% =============================================================================
% Chapter 2: Environment Setup
% =============================================================================
\chapter{Environment Setup}
\label{ch:environment-setup}

\begin{objectives}
    \item Install Python and the \texttt{uv} package manager
    \item Clone and set up the TryLM repository
    \item Understand the project structure
    \item Run basic CLI commands
    \item (Optional) Configure GPU acceleration
\end{objectives}

% -----------------------------------------------------------------------------
\section{Prerequisites}
% -----------------------------------------------------------------------------

Before we begin, ensure you have:

\begin{itemize}
    \item A computer running Linux, macOS, or Windows (with WSL2)
    \item Python 3.11 or later
    \item Basic familiarity with the command line
    \item (Optional) An NVIDIA GPU with CUDA support for faster training
\end{itemize}

% -----------------------------------------------------------------------------
\section{Installing the uv Package Manager}
% -----------------------------------------------------------------------------

\TryLM uses \texttt{uv}, a modern, fast Python package manager written in Rust. It replaces traditional tools like \texttt{pip} and \texttt{virtualenv} with a single, unified tool.

\subsection{Why uv?}

\begin{itemize}
    \item \textbf{Speed}: 10--100x faster than pip
    \item \textbf{Reliability}: Deterministic dependency resolution
    \item \textbf{Simplicity}: Single tool for environments and packages
\end{itemize}

\subsection{Installation}

Install \texttt{uv} using the official installer:

\begin{lstlisting}[style=shell]
# Linux/macOS
curl -LsSf https://astral.sh/uv/install.sh | sh

# Windows (PowerShell)
powershell -c "irm https://astral.sh/uv/install.ps1 | iex"
\end{lstlisting}

After installation, restart your terminal and verify:

\begin{lstlisting}[style=shell]
uv --version
# uv 0.4.x
\end{lstlisting}

% -----------------------------------------------------------------------------
\section{Getting the TryLM Code}
% -----------------------------------------------------------------------------

\subsection{Clone the Repository}

\begin{lstlisting}[style=shell]
git clone https://github.com/your-repo/trylm.git
cd trylm
\end{lstlisting}

\subsection{Install Dependencies}

Use \texttt{uv} to create a virtual environment and install all dependencies:

\begin{lstlisting}[style=shell]
uv sync
\end{lstlisting}

This command:
\begin{enumerate}
    \item Creates a \texttt{.venv} virtual environment
    \item Reads \texttt{pyproject.toml} for dependency specifications
    \item Installs all required packages
\end{enumerate}

\begin{note}
All \TryLM commands use \texttt{uv run} to execute within the virtual environment. You don't need to manually activate the environment.
\end{note}

% -----------------------------------------------------------------------------
\section{Project Structure}
% -----------------------------------------------------------------------------

Let's explore the \TryLM codebase:

\begin{lstlisting}[style=shell, numbers=none]
trylm/
|-- pyproject.toml        # Project configuration and dependencies
|-- configs/
|   `-- tiny_stories.yaml # Default training configuration
|-- src/
|   `-- trylm/
|       |-- __init__.py   # Package init
|       |-- cli.py        # Command-line interface
|       |-- config.py     # Configuration dataclasses
|       |-- model.py      # GPT model architecture
|       |-- tokenizer.py  # BPE tokenizer
|       |-- data.py       # Dataset and dataloaders
|       |-- trainer.py    # Training loop
|       |-- generate.py   # Text generation
|       `-- evaluate.py   # Evaluation metrics
|-- data/                 # Tokenizer and cached data
|-- checkpoints/          # Saved model checkpoints
`-- notebooks/            # Jupyter notebooks
\end{lstlisting}

\subsection{Key Files}

\begin{table}[H]
\centering
\begin{tabular}{lp{8cm}}
\toprule
\textbf{File} & \textbf{Purpose} \\
\midrule
\texttt{model.py} & The GPT model: embeddings, attention, feed-forward layers, and the complete architecture \\
\texttt{tokenizer.py} & BPE tokenizer training and text encoding/decoding \\
\texttt{data.py} & Dataset loading, tokenization, and batch preparation \\
\texttt{trainer.py} & Training loop with optimization, checkpointing, and logging \\
\texttt{generate.py} & Text generation with various sampling strategies \\
\texttt{config.py} & Configuration dataclasses for tokenizer, model, and training \\
\texttt{cli.py} & Typer-based command-line interface \\
\bottomrule
\end{tabular}
\caption{Key source files in the TryLM project}
\end{table}

% -----------------------------------------------------------------------------
\section{The Command-Line Interface}
% -----------------------------------------------------------------------------

\TryLM provides a command-line interface (CLI) for all operations. The entry point is the \texttt{trylm} command.

\subsection{Available Commands}

\begin{lstlisting}[style=shell]
uv run trylm --help
\end{lstlisting}

This shows:
\begin{lstlisting}[style=shell, numbers=none]
Usage: trylm [OPTIONS] COMMAND [ARGS]...

Commands:
  tokenizer  Tokenizer training and encoding
  train      Train the language model
  generate   Generate text from a trained model
  eval       Evaluate model perplexity
  info       Display model information
  init       Initialize a configuration file
\end{lstlisting}

\subsection{Your First Command}

Let's check the default model configuration:

\begin{lstlisting}[style=shell]
uv run trylm info
\end{lstlisting}

Output:
\begin{lstlisting}[style=shell, numbers=none]
TryLM Model Configuration
==========================
Vocabulary size:  16,384
Context length:   512
Model dimension:  384
Attention heads:  6
Layers:           6
FFN dimension:    1,536

Estimated parameters: ~10.8M
\end{lstlisting}

\begin{keypoint}
The \texttt{trylm info} command shows the model architecture without loading any weights. It's useful for understanding the configuration before training.
\end{keypoint}

% -----------------------------------------------------------------------------
\section{Configuration System}
% -----------------------------------------------------------------------------

\TryLM uses a hierarchical configuration system defined in \texttt{config.py}.

\subsection{Configuration Dataclasses}

Three dataclasses define the complete configuration:

\begin{lstlisting}[style=python, caption={Configuration structure from config.py}]
@dataclass
class TokenizerConfig:
    vocab_size: int = 16384
    min_frequency: int = 2
    special_tokens: list[str] = field(default_factory=lambda: [
        "<|pad|>", "<|eos|>", "<|bos|>", "<|unk|>"
    ])
    save_path: Path = Path("data/tokenizer")

@dataclass
class ModelConfig:
    vocab_size: int = 16384
    n_layers: int = 6
    n_heads: int = 6
    d_model: int = 384
    d_ff: int = 1536
    context_length: int = 512
    dropout: float = 0.1
    bias: bool = False

@dataclass
class TrainConfig:
    batch_size: int = 64
    learning_rate: float = 3e-4
    max_steps: int = 10000
    # ... more options
\end{lstlisting}

\subsection{YAML Configuration Files}

Configuration can be loaded from YAML files. Here's the default \texttt{tiny\_stories.yaml}:

\begin{lstlisting}[style=yaml, caption={configs/tiny\_stories.yaml}]
tokenizer:
  vocab_size: 16384
  min_frequency: 2

model:
  vocab_size: 16384
  n_layers: 6
  n_heads: 6
  d_model: 384
  d_ff: 1536
  context_length: 512
  dropout: 0.1
  bias: false

train:
  dataset_name: roneneldan/TinyStories
  batch_size: 64
  gradient_accumulation_steps: 4
  learning_rate: 0.0003
  max_steps: 10000
  warmup_steps: 100
  eval_interval: 500
  checkpoint_dir: checkpoints
  mixed_precision: true
  compile: true
  seed: 42
\end{lstlisting}

\subsection{CLI Overrides}

Any configuration value can be overridden via command-line arguments:

\begin{lstlisting}[style=shell]
# Use config file with overrides
uv run trylm train --config configs/tiny_stories.yaml \
    --batch-size 32 \
    --max-steps 5000
\end{lstlisting}

% -----------------------------------------------------------------------------
\section{Quick Test: Training the Tokenizer}
% -----------------------------------------------------------------------------

Let's verify everything works by training a tokenizer:

\begin{lstlisting}[style=shell]
uv run trylm tokenizer train --vocab-size 16384
\end{lstlisting}

This command:
\begin{enumerate}
    \item Downloads the TinyStories dataset (first time only)
    \item Trains a BPE tokenizer on the text
    \item Saves the tokenizer to \texttt{data/tokenizer/}
\end{enumerate}

Expected output:
\begin{lstlisting}[style=shell, numbers=none]
Training tokenizer on roneneldan/TinyStories...
Processing 2,119,719 examples...
Training BPE tokenizer with vocab_size=16384...
Tokenizer saved to data/tokenizer
\end{lstlisting}

Now test encoding some text:

\begin{lstlisting}[style=shell]
uv run trylm tokenizer encode "Once upon a time"
\end{lstlisting}

Output:
\begin{lstlisting}[style=shell, numbers=none]
Text: Once upon a time
Token IDs: [1, 4521, 2378, 263, 1159, 2]
Tokens: ['<|bos|>', 'Once', ' upon', ' a', ' time', '<|eos|>']
\end{lstlisting}

% -----------------------------------------------------------------------------
\section{GPU Setup (Optional)}
% -----------------------------------------------------------------------------

Training is significantly faster on a GPU. \TryLM automatically detects and uses available GPUs.

\subsection{NVIDIA CUDA}

For NVIDIA GPUs, ensure you have:
\begin{enumerate}
    \item NVIDIA driver installed
    \item CUDA toolkit 11.8 or later
    \item cuDNN library
\end{enumerate}

Verify CUDA availability in Python:

\begin{lstlisting}[style=python]
import torch
print(torch.cuda.is_available())  # True
print(torch.cuda.get_device_name(0))  # e.g., "NVIDIA GeForce RTX 3080"
\end{lstlisting}

\subsection{Apple Silicon (MPS)}

On Apple Silicon Macs (M1/M2/M3), PyTorch uses the Metal Performance Shaders (MPS) backend:

\begin{lstlisting}[style=python]
import torch
print(torch.backends.mps.is_available())  # True on Apple Silicon
\end{lstlisting}

\subsection{Device Selection}

\TryLM automatically selects the best available device:

\begin{lstlisting}[style=python, caption={Device selection from cli.py}]
def get_device() -> str:
    if torch.cuda.is_available():
        return "cuda"
    elif torch.backends.mps.is_available():
        return "mps"
    return "cpu"
\end{lstlisting}

\begin{note}
Training on CPU is possible but slow. A 10M parameter model takes:
\begin{itemize}
    \item GPU (RTX 3080): $\sim$1 hour for 10k steps
    \item CPU (modern): $\sim$10+ hours for 10k steps
\end{itemize}
\end{note}

% -----------------------------------------------------------------------------
\section{Weights \& Biases Setup (Optional)}
% -----------------------------------------------------------------------------

\TryLM integrates with Weights \& Biases (W\&B) for experiment tracking.

\subsection{Create an Account}

Sign up at \url{https://wandb.ai} (free for personal use).

\subsection{Login}

\begin{lstlisting}[style=shell]
uv run wandb login
\end{lstlisting}

This prompts you to enter your API key.

\subsection{Verify}

When you train with W\&B enabled, you'll see:
\begin{lstlisting}[style=shell, numbers=none]
wandb: Syncing run trylm-tiny-stories
wandb: View project at https://wandb.ai/your-username/trylm
\end{lstlisting}

W\&B provides:
\begin{itemize}
    \item Real-time loss and metric plots
    \item Hyperparameter tracking
    \item Model checkpoint versioning
    \item Experiment comparison
\end{itemize}

% -----------------------------------------------------------------------------
\section{Development Tools}
% -----------------------------------------------------------------------------

For exploring and modifying the code, we recommend:

\subsection{IDE/Editor}

\begin{itemize}
    \item \textbf{VS Code} with Python extension
    \item \textbf{PyCharm} (Community or Professional)
    \item \textbf{Neovim} with LSP support
\end{itemize}

\subsection{Useful Commands}

\begin{lstlisting}[style=shell]
# Run type checking
uv run mypy src/trylm

# Run tests (if available)
uv run pytest

# Format code
uv run ruff format src/

# Lint code
uv run ruff check src/
\end{lstlisting}

% -----------------------------------------------------------------------------
\section{Troubleshooting}
% -----------------------------------------------------------------------------

\begin{pitfall}[title=Common Issues]
\textbf{``Command not found: uv''}

Ensure \texttt{uv} is in your PATH. Restart your terminal after installation.

\textbf{``CUDA out of memory''}

Reduce batch size: \texttt{--batch-size 16} or disable mixed precision.

\textbf{``Module not found: trylm''}

Run \texttt{uv sync} to install the package in development mode.

\textbf{Slow download of TinyStories}

The dataset is $\sim$500MB. First download may take several minutes.
\end{pitfall}

% -----------------------------------------------------------------------------
\section{Exercises}
% -----------------------------------------------------------------------------

\begin{exercise}
\textbf{Explore the CLI}

Run each of these commands and observe the output:
\begin{lstlisting}[style=shell, numbers=none]
uv run trylm --help
uv run trylm tokenizer --help
uv run trylm train --help
uv run trylm generate --help
\end{lstlisting}
What options are available for each command?
\end{exercise}

\begin{exercise}
\textbf{Examine the Project Structure}

Open \texttt{src/trylm/model.py} and identify:
\begin{enumerate}
    \item The main model class
    \item The attention mechanism class
    \item The feed-forward network class
\end{enumerate}
Don't worry about understanding the code yet---just get familiar with the structure.
\end{exercise}

\begin{exercise}
\textbf{Configuration Exploration}

Create a custom configuration file:
\begin{lstlisting}[style=shell, numbers=none]
uv run trylm init --output my_config.yaml
\end{lstlisting}
Open the file and identify:
\begin{enumerate}
    \item How many transformer layers are configured?
    \item What is the vocabulary size?
    \item What is the learning rate?
\end{enumerate}
\end{exercise}

% -----------------------------------------------------------------------------
\section{Summary}
% -----------------------------------------------------------------------------

\begin{summary}
    \item \texttt{uv} is a fast, modern Python package manager
    \item All commands use \texttt{uv run trylm <command>}
    \item The project follows a clean \texttt{src/} layout with separate modules for each concern
    \item Configuration uses dataclasses and YAML files
    \item CLI arguments can override configuration values
    \item GPU acceleration is automatic when available
    \item W\&B provides optional experiment tracking
\end{summary}

In the next chapter, we'll dive into tokenization---the critical first step in processing text for language models.
